\chapter{\abstractname}

Multilingual Large Language Models (LLMs) demonstrate impressive zero-shot cross-lingual capabilities, yet the mechanisms governing how internal representations align across languages remain incompletely understood. This thesis systematically evaluates cross-lingual alignment in modern LLMs using the MEXA (Multilingual Evaluation via Cross-Lingual Alignment) framework. We reproduce reference evaluations on standard decoders and extend the study along four major axes: contemporary decoder families (Qwen 3, Qwen 3.5, Apertus 1.5, Gemma 4), sparse Mixture-of-Experts (MoE) architectures (Mixtral 8x7B, Mixtral 8x22B), specialized sentence encoders, and non-English pivot languages (French, German, Arabic, Chinese, Basque).

Our empirical analysis across FLORES-200 and the Parallel Bible Corpus yields four key insights: (1) Parameter scaling and multilingual training recipes substantially enhance peak representation alignment, with Qwen 3.5 9B ($\mu_{\text{Max}} = 0.7794$) and Swiss AI Apertus 1.5 70B setting state-of-the-art decoder scores ($\mu_{\text{Max}} = 0.5036$ on Bible); (2) Dynamic token-level routing in sparse MoE models boosts cross-lingual alignment over dense counterparts without fragmenting the shared interlingua; (3) Substituting English for non-English pivot anchors retains high retrieval scores, and mean-centering eliminates language-identity centroid offsets, achieving uniform peak alignment ($\mu_{\text{Max}} \ge 0.89 - 0.95$); and (4) MEXA alignment scores exhibit strong predictive validity against downstream multilingual benchmarks including Belebele ($\rho = 0.8420$) and m-ARC ($\rho = 0.8105$). These findings validate MEXA as a robust, computationally efficient proxy for evaluating multilingual LLM capabilities.
